% --- Archon physics formalization source begin ---
% archon:physics
% archon:covers PhyXMiniProblems/problem_phyx_mini_0386.lean
% archon:source-report reports/phyx_mini/problem_phyx_mini_0386.source.json

\chapter{Physics problem phyx\_mini\_0386}
\label{ch:PhyXMiniProblems_problem_phyx_mini_0386}

\paragraph{Problem source.}
\#\# Physical scenario

A steel tank of cross-sectional area 3 \textbackslash{}m\textasciicircum{}2\textbackslash{} and height 16 m weighs 10 000 kg and is open at the top, as shown in the figure. We want to float it in the ocean so that it is positioned 10 m straight down by pouring concrete into its bottom.

\#\# Question

How much concrete should we use?

\#\# Answer choices

- A: A: 490kg

- B: B: 19910kg

- C: C: 154kg

- D: D: 10.2kg

\#\# Figure caption (auxiliary; use the image as primary evidence)

The image shows a diagram of a vertical structure partially submerged in the ocean. The structure is a U-shaped container with the following components:

1. **Air**: The top part of the container is filled with air.
2. **Concrete**: Below the air, a segment is filled with concrete, illustrated in blue.
3. **Ocean**: The container is in the ocean, depicted by a wavy line representing the water surface.
4. **Depth Indication**: A vertical arrow labeled "10 m" indicates the distance from the ocean surface to the bottom of the concrete section.
5. **Labels**:
   - The word "Air" is placed within the container's top section.
   - "Concrete" is labeled within the blue section.
   - "Ocean" is labeled outside the container, parallel to the water surface. 

These elements collectively illustrate a submerged concrete-filled container with specified sections and measurements.

\#\# Dataset metadata

Laws of Thermodynamics; Physical Model Grounding Reasoning, Multi-Formula Reasoning

\paragraph{Recorded answer/context.}
B: B: 19910kg

\paragraph{Figure/image path.}
/root/proposal\_for\_physic/science-mango/phyx\_mini\_run/phyx\_data/test\_image/386.png

\paragraph{Formalization target.}
create a compiling Lean file with sorry bodies at `PhyXMiniProblems/problem\_phyx\_mini\_0386.lean`. The Lean declarations must preserve the physical quantities, dimensions or dimensional roles, figure labels, governing-law hypotheses, and final relation expressed by this problem.
Use Mathlib/Physlib names found through LeanExplore where available. If a physics API is missing, introduce faithful local abstractions rather than scalar placeholder aliases.

\begin{theorem}[Physics formalization target]
\label{thm:physics:phyx_mini_0386:target}
\lean{PhyXMiniProblems.ProblemPhyXMini0386.requiredConcreteMass_eq_answerB}
\uses{def:physics:phyx-mini-0386:phyxminiproblems-problemphyxmini0386-kilograms, def:physics:phyx-mini-0386:phyxminiproblems-problemphyxmini0386-ballastedtanksetup, def:physics:phyx-mini-0386:phyxminiproblems-problemphyxmini0386-hastankandfiguredata, def:physics:phyx-mini-0386:phyxminiproblems-problemphyxmini0386-hasrecordedoceandensitycalibration, def:physics:phyx-mini-0386:phyxminiproblems-problemphyxmini0386-hasphysicalgravity, def:physics:phyx-mini-0386:phyxminiproblems-problemphyxmini0386-satisfiesprismaticdisplacementlaw, def:physics:phyx-mini-0386:phyxminiproblems-problemphyxmini0386-satisfiestotalweightlaw, def:physics:phyx-mini-0386:phyxminiproblems-problemphyxmini0386-satisfiesarchimedeslaw, def:physics:phyx-mini-0386:phyxminiproblems-problemphyxmini0386-floatsinverticalstaticequilibrium, def:physics:phyx-mini-0386:phyxminiproblems-problemphyxmini0386-answermassinkilograms}
The assigned autoformalize agent should translate this physics problem into Lean declarations in the covered file, with theorem and lemma proof bodies written as `by sorry`.
\end{theorem}
\begin{proof}
This is an autoformalization task, not a proof task. Produce statements that can later be proved by the physics prover without weakening the physical model.
\end{proof}

% --- Archon declaration topology begin ---
\section{Declaration topology}
\begin{definition}[area Dimension]
  \label{def:physics:phyx-mini-0386:phyxminiproblems-problemphyxmini0386-areadimension}
  \lean{PhyXMiniProblems.ProblemPhyXMini0386.areaDimension}
  The physical dimension of an area
\end{definition}
\begin{proof}
  This is the declared model definition of area Dimension.
\end{proof}
\begin{definition}[volume Dimension]
  \label{def:physics:phyx-mini-0386:phyxminiproblems-problemphyxmini0386-volumedimension}
  \lean{PhyXMiniProblems.ProblemPhyXMini0386.volumeDimension}
  The physical dimension of a volume
\end{definition}
\begin{proof}
  This is the declared model definition of volume Dimension.
\end{proof}
\begin{definition}[mass Density Dimension]
  \label{def:physics:phyx-mini-0386:phyxminiproblems-problemphyxmini0386-massdensitydimension}
  \lean{PhyXMiniProblems.ProblemPhyXMini0386.massDensityDimension}
  The physical dimension of mass density, M L⁻³
\end{definition}
\begin{proof}
  This is the declared model definition of mass Density Dimension.
\end{proof}
\begin{definition}[acceleration Dimension]
  \label{def:physics:phyx-mini-0386:phyxminiproblems-problemphyxmini0386-accelerationdimension}
  \lean{PhyXMiniProblems.ProblemPhyXMini0386.accelerationDimension}
  The physical dimension of acceleration
\end{definition}
\begin{proof}
  This is the declared model definition of acceleration Dimension.
\end{proof}
\begin{definition}[force Dimension]
  \label{def:physics:phyx-mini-0386:phyxminiproblems-problemphyxmini0386-forcedimension}
  \lean{PhyXMiniProblems.ProblemPhyXMini0386.forceDimension}
  The physical dimension of force
\end{definition}
\begin{proof}
  This is the declared model definition of force Dimension.
\end{proof}
\begin{definition}[square Meters]
  \label{def:physics:phyx-mini-0386:phyxminiproblems-problemphyxmini0386-squaremeters}
  \lean{PhyXMiniProblems.ProblemPhyXMini0386.squareMeters}
  \uses{def:physics:phyx-mini-0386:phyxminiproblems-problemphyxmini0386-areadimension}
  An area readout in square metres in the SI unit system
\end{definition}
\begin{proof}
  This is the declared model definition of square Meters.
\end{proof}
\begin{definition}[meters]
  \label{def:physics:phyx-mini-0386:phyxminiproblems-problemphyxmini0386-meters}
  \lean{PhyXMiniProblems.ProblemPhyXMini0386.meters}
  A length readout in metres in the SI unit system
\end{definition}
\begin{proof}
  This is the declared model definition of meters.
\end{proof}
\begin{definition}[cubic Meters]
  \label{def:physics:phyx-mini-0386:phyxminiproblems-problemphyxmini0386-cubicmeters}
  \lean{PhyXMiniProblems.ProblemPhyXMini0386.cubicMeters}
  \uses{def:physics:phyx-mini-0386:phyxminiproblems-problemphyxmini0386-volumedimension}
  A volume readout in cubic metres in the SI unit system
\end{definition}
\begin{proof}
  This is the declared model definition of cubic Meters.
\end{proof}
\begin{definition}[kilograms]
  \label{def:physics:phyx-mini-0386:phyxminiproblems-problemphyxmini0386-kilograms}
  \lean{PhyXMiniProblems.ProblemPhyXMini0386.kilograms}
  A mass readout in kilograms in the SI unit system
\end{definition}
\begin{proof}
  This is the declared model definition of kilograms.
\end{proof}
\begin{definition}[kilograms Per Cubic Meter]
  \label{def:physics:phyx-mini-0386:phyxminiproblems-problemphyxmini0386-kilogramspercubicmeter}
  \lean{PhyXMiniProblems.ProblemPhyXMini0386.kilogramsPerCubicMeter}
  \uses{def:physics:phyx-mini-0386:phyxminiproblems-problemphyxmini0386-massdensitydimension}
  A mass-density readout in kilograms per cubic metre
\end{definition}
\begin{proof}
  This is the declared model definition of kilograms Per Cubic Meter.
\end{proof}
\begin{definition}[meters Per Second Squared]
  \label{def:physics:phyx-mini-0386:phyxminiproblems-problemphyxmini0386-meterspersecondsquared}
  \lean{PhyXMiniProblems.ProblemPhyXMini0386.metersPerSecondSquared}
  \uses{def:physics:phyx-mini-0386:phyxminiproblems-problemphyxmini0386-accelerationdimension}
  An acceleration readout in metres per second squared
\end{definition}
\begin{proof}
  This is the declared model definition of meters Per Second Squared.
\end{proof}
\begin{definition}[newtons]
  \label{def:physics:phyx-mini-0386:phyxminiproblems-problemphyxmini0386-newtons}
  \lean{PhyXMiniProblems.ProblemPhyXMini0386.newtons}
  \uses{def:physics:phyx-mini-0386:phyxminiproblems-problemphyxmini0386-forcedimension}
  A force readout in newtons
\end{definition}
\begin{proof}
  This is the declared model definition of newtons.
\end{proof}
\begin{definition}[Tank Top Boundary]
  \label{def:physics:phyx-mini-0386:phyxminiproblems-problemphyxmini0386-tanktopboundary}
  \lean{PhyXMiniProblems.ProblemPhyXMini0386.TankTopBoundary}
  Whether the tank's upper boundary is open or sealed
\end{definition}
\begin{proof}
  This is the declared model definition of Tank Top Boundary.
\end{proof}
\begin{definition}[Tank Axis Orientation]
  \label{def:physics:phyx-mini-0386:phyxminiproblems-problemphyxmini0386-tankaxisorientation}
  \lean{PhyXMiniProblems.ProblemPhyXMini0386.TankAxisOrientation}
  Orientation of the tank's long axis relative to the ocean surface
\end{definition}
\begin{proof}
  This is the declared model definition of Tank Axis Orientation.
\end{proof}
\begin{definition}[Figure Material]
  \label{def:physics:phyx-mini-0386:phyxminiproblems-problemphyxmini0386-figurematerial}
  \lean{PhyXMiniProblems.ProblemPhyXMini0386.FigureMaterial}
  Material names stated in the scenario or printed in the figure
\end{definition}
\begin{proof}
  This is the declared model definition of Figure Material.
\end{proof}
\begin{definition}[Ballast Location]
  \label{def:physics:phyx-mini-0386:phyxminiproblems-problemphyxmini0386-ballastlocation}
  \lean{PhyXMiniProblems.ProblemPhyXMini0386.BallastLocation}
  The location of the concrete ballast inside the tank
\end{definition}
\begin{proof}
  This is the declared model definition of Ballast Location.
\end{proof}
\begin{definition}[Ballasted Tank Setup]
  \label{def:physics:phyx-mini-0386:phyxminiproblems-problemphyxmini0386-ballastedtanksetup}
  \lean{PhyXMiniProblems.ProblemPhyXMini0386.BallastedTankSetup}
  \uses{def:physics:phyx-mini-0386:phyxminiproblems-problemphyxmini0386-areadimension, def:physics:phyx-mini-0386:phyxminiproblems-problemphyxmini0386-volumedimension, def:physics:phyx-mini-0386:phyxminiproblems-problemphyxmini0386-massdensitydimension, def:physics:phyx-mini-0386:phyxminiproblems-problemphyxmini0386-accelerationdimension, def:physics:phyx-mini-0386:phyxminiproblems-problemphyxmini0386-forcedimension, def:physics:phyx-mini-0386:phyxminiproblems-problemphyxmini0386-tanktopboundary, def:physics:phyx-mini-0386:phyxminiproblems-problemphyxmini0386-tankaxisorientation, def:physics:phyx-mini-0386:phyxminiproblems-problemphyxmini0386-figurematerial, def:physics:phyx-mini-0386:phyxminiproblems-problemphyxmini0386-ballastlocation}
  Physical quantities and categorical figure data for the ballasted tank. The concrete mass and the force magnitudes are independent observables. Their relations are supplied only by the governing-law predicates below, so the requested mass is not built into the setup
\end{definition}
\begin{proof}
  This is the declared model definition of Ballasted Tank Setup.
\end{proof}
\begin{definition}[total Supported Mass]
  \label{def:physics:phyx-mini-0386:phyxminiproblems-problemphyxmini0386-totalsupportedmass}
  \lean{PhyXMiniProblems.ProblemPhyXMini0386.totalSupportedMass}
  \uses{def:physics:phyx-mini-0386:phyxminiproblems-problemphyxmini0386-ballastedtanksetup}
  The total supported mass is the steel tank mass plus its concrete ballast. The mass of the contained air is neglected in this elementary model
\end{definition}
\begin{proof}
  This is the declared model definition of total Supported Mass.
\end{proof}
\begin{definition}[Has Tank And Figure Data]
  \label{def:physics:phyx-mini-0386:phyxminiproblems-problemphyxmini0386-hastankandfiguredata}
  \lean{PhyXMiniProblems.ProblemPhyXMini0386.HasTankAndFigureData}
  \uses{def:physics:phyx-mini-0386:phyxminiproblems-problemphyxmini0386-squaremeters, def:physics:phyx-mini-0386:phyxminiproblems-problemphyxmini0386-meters, def:physics:phyx-mini-0386:phyxminiproblems-problemphyxmini0386-kilograms, def:physics:phyx-mini-0386:phyxminiproblems-problemphyxmini0386-ballastedtanksetup}
  Numerical and diagrammatic data explicitly supplied by the problem and its figure. In particular, the 10 m arrow is read as the depth of the tank bottom below the ocean surface, not as the concrete fill height
\end{definition}
\begin{proof}
  This is the declared model definition of Has Tank And Figure Data.
\end{proof}
\begin{definition}[Has Recorded Ocean Density Calibration]
  \label{def:physics:phyx-mini-0386:phyxminiproblems-problemphyxmini0386-hasrecordedoceandensitycalibration}
  \lean{PhyXMiniProblems.ProblemPhyXMini0386.HasRecordedOceanDensityCalibration}
  \uses{def:physics:phyx-mini-0386:phyxminiproblems-problemphyxmini0386-kilogramspercubicmeter, def:physics:phyx-mini-0386:phyxminiproblems-problemphyxmini0386-ballastedtanksetup}
  Reference-density calibration implicit in the recorded exact answer. The problem does not state an ocean density. The choice 19910 kg is obtained from the textbook readout 997 kg/m³; keeping this as a separate hypothesis makes that otherwise hidden modeling choice explicit
\end{definition}
\begin{proof}
  This is the declared model definition of Has Recorded Ocean Density Calibration.
\end{proof}
\begin{definition}[Has Physical Gravity]
  \label{def:physics:phyx-mini-0386:phyxminiproblems-problemphyxmini0386-hasphysicalgravity}
  \lean{PhyXMiniProblems.ProblemPhyXMini0386.HasPhysicalGravity}
  \uses{def:physics:phyx-mini-0386:phyxminiproblems-problemphyxmini0386-ballastedtanksetup}
  Nondegeneracy needed to cancel gravity from the force balance
\end{definition}
\begin{proof}
  This is the declared model definition of Has Physical Gravity.
\end{proof}
\begin{definition}[Satisfies Prismatic Displacement Law]
  \label{def:physics:phyx-mini-0386:phyxminiproblems-problemphyxmini0386-satisfiesprismaticdisplacementlaw}
  \lean{PhyXMiniProblems.ProblemPhyXMini0386.SatisfiesPrismaticDisplacementLaw}
  \uses{def:physics:phyx-mini-0386:phyxminiproblems-problemphyxmini0386-areadimension, def:physics:phyx-mini-0386:phyxminiproblems-problemphyxmini0386-volumedimension, def:physics:phyx-mini-0386:phyxminiproblems-problemphyxmini0386-ballastedtanksetup}
  Prismatic displacement law for the upright, unflooded tank: the displaced ocean volume is the effective horizontal cross-sectional area times the draft. Steel-wall thickness is neglected at the resolution of the supplied area
\end{definition}
\begin{proof}
  This is the declared model definition of Satisfies Prismatic Displacement Law.
\end{proof}
\begin{definition}[Satisfies Total Weight Law]
  \label{def:physics:phyx-mini-0386:phyxminiproblems-problemphyxmini0386-satisfiestotalweightlaw}
  \lean{PhyXMiniProblems.ProblemPhyXMini0386.SatisfiesTotalWeightLaw}
  \uses{def:physics:phyx-mini-0386:phyxminiproblems-problemphyxmini0386-accelerationdimension, def:physics:phyx-mini-0386:phyxminiproblems-problemphyxmini0386-forcedimension, def:physics:phyx-mini-0386:phyxminiproblems-problemphyxmini0386-ballastedtanksetup, def:physics:phyx-mini-0386:phyxminiproblems-problemphyxmini0386-totalsupportedmass}
  The tank's downward weight magnitude is (tank mass + ballast mass) g
\end{definition}
\begin{proof}
  This is the declared model definition of Satisfies Total Weight Law.
\end{proof}
\begin{definition}[Satisfies Archimedes Law]
  \label{def:physics:phyx-mini-0386:phyxminiproblems-problemphyxmini0386-satisfiesarchimedeslaw}
  \lean{PhyXMiniProblems.ProblemPhyXMini0386.SatisfiesArchimedesLaw}
  \uses{def:physics:phyx-mini-0386:phyxminiproblems-problemphyxmini0386-volumedimension, def:physics:phyx-mini-0386:phyxminiproblems-problemphyxmini0386-massdensitydimension, def:physics:phyx-mini-0386:phyxminiproblems-problemphyxmini0386-accelerationdimension, def:physics:phyx-mini-0386:phyxminiproblems-problemphyxmini0386-forcedimension, def:physics:phyx-mini-0386:phyxminiproblems-problemphyxmini0386-ballastedtanksetup}
  Archimedes' law: buoyancy is ocean density × displaced volume × g
\end{definition}
\begin{proof}
  This is the declared model definition of Satisfies Archimedes Law.
\end{proof}
\begin{definition}[Floats In Vertical Static Equilibrium]
  \label{def:physics:phyx-mini-0386:phyxminiproblems-problemphyxmini0386-floatsinverticalstaticequilibrium}
  \lean{PhyXMiniProblems.ProblemPhyXMini0386.FloatsInVerticalStaticEquilibrium}
  \uses{def:physics:phyx-mini-0386:phyxminiproblems-problemphyxmini0386-ballastedtanksetup}
  Static vertical flotation: the upward buoyant force balances the weight
\end{definition}
\begin{proof}
  This is the declared model definition of Floats In Vertical Static Equilibrium.
\end{proof}
\begin{definition}[Answer Choice]
  \label{def:physics:phyx-mini-0386:phyxminiproblems-problemphyxmini0386-answerchoice}
  \lean{PhyXMiniProblems.ProblemPhyXMini0386.AnswerChoice}
  The four concrete-mass answers printed beside the multiple-choice labels
\end{definition}
\begin{proof}
  This is the declared model definition of Answer Choice.
\end{proof}
\begin{definition}[answer Mass In Kilograms]
  \label{def:physics:phyx-mini-0386:phyxminiproblems-problemphyxmini0386-answermassinkilograms}
  \lean{PhyXMiniProblems.ProblemPhyXMini0386.answerMassInKilograms}
  \uses{def:physics:phyx-mini-0386:phyxminiproblems-problemphyxmini0386-answerchoice}
  The concrete mass, in kilograms, displayed for an answer choice
\end{definition}
\begin{proof}
  This is the declared model definition of answer Mass In Kilograms.
\end{proof}
\begin{lemma}[displaced Ocean Volume at requested Draft]
  \label{lem:physics:phyx-mini-0386:phyxminiproblems-problemphyxmini0386-displacedoceanvolume-at-requesteddraft}
  \lean{PhyXMiniProblems.ProblemPhyXMini0386.displacedOceanVolume_at_requestedDraft}
  \uses{def:physics:phyx-mini-0386:phyxminiproblems-problemphyxmini0386-cubicmeters, def:physics:phyx-mini-0386:phyxminiproblems-problemphyxmini0386-ballastedtanksetup, def:physics:phyx-mini-0386:phyxminiproblems-problemphyxmini0386-hastankandfiguredata, def:physics:phyx-mini-0386:phyxminiproblems-problemphyxmini0386-satisfiesprismaticdisplacementlaw}
  The area and requested 10 m draft displace exactly 30 m³ of ocean
\end{lemma}
\begin{proof}
  The conclusion follows from the governing relations and figure readouts named in the statement of displaced Ocean Volume at requested Draft.
\end{proof}
\begin{lemma}[total Supported Mass at requested Draft]
  \label{lem:physics:phyx-mini-0386:phyxminiproblems-problemphyxmini0386-totalsupportedmass-at-requesteddraft}
  \lean{PhyXMiniProblems.ProblemPhyXMini0386.totalSupportedMass_at_requestedDraft}
  \uses{def:physics:phyx-mini-0386:phyxminiproblems-problemphyxmini0386-kilograms, def:physics:phyx-mini-0386:phyxminiproblems-problemphyxmini0386-ballastedtanksetup, def:physics:phyx-mini-0386:phyxminiproblems-problemphyxmini0386-totalsupportedmass, def:physics:phyx-mini-0386:phyxminiproblems-problemphyxmini0386-hastankandfiguredata, def:physics:phyx-mini-0386:phyxminiproblems-problemphyxmini0386-hasrecordedoceandensitycalibration, def:physics:phyx-mini-0386:phyxminiproblems-problemphyxmini0386-hasphysicalgravity, def:physics:phyx-mini-0386:phyxminiproblems-problemphyxmini0386-satisfiesprismaticdisplacementlaw, def:physics:phyx-mini-0386:phyxminiproblems-problemphyxmini0386-satisfiestotalweightlaw, def:physics:phyx-mini-0386:phyxminiproblems-problemphyxmini0386-satisfiesarchimedeslaw, def:physics:phyx-mini-0386:phyxminiproblems-problemphyxmini0386-floatsinverticalstaticequilibrium}
  At the requested draft, buoyancy supports 997 × 30 = 29910 kg
\end{lemma}
\begin{proof}
  The conclusion follows from the governing relations and figure readouts named in the statement of total Supported Mass at requested Draft.
\end{proof}
% --- Archon declaration topology end ---
% --- Archon physics formalization source end ---
